% Notebook 14 body fragment; no preamble. Unique prefix nw-.
% Requires amsmath, amssymb, amsthm and \norm; follows nc-/nb- fragments.
% Main proof independently audited; source action convention scope explicit.

\section{A nonflat Wilson Hessian and its full energy inverse}
\label{nw-section}

The averaging bounds above close the remaining averaging hypotheses
of an energy-form perturbation argument. We first calculate the
Wilson Hessian, including its normalization. The resulting theorem
is for an explicitly defined Wilson operator with the deposited
homogeneous averaging maps. Identification with every reused source
symbol for an action Hessian is stated separately.

\subsection{Action, charts, and an exact local decomposition}
\label{nw-wilson}

Write $\eta=M^{-1}\le1$ and $\tau=N_c^{-1}\operatorname{Tr}$.
Fine fields have physical pairing
$\langle B,C\rangle_\eta=\eta^d\sum_e\tau(B_eC_e)$.
The sum is over positive bonds; plaquettes below are counted once,
with $\mu<\nu$. On Hermitian fields this is the real Hilbert pairing;
its complex extension is bilinear. Superscript $T$ always denotes
that transpose, with no conjugation of a complex background.
Define
\begin{equation}\label{nw-action}
 \mathcal S_\eta(U)=\eta^{d-4}\sum_p
   \left(1-\frac12\tau(P_p+P_p^{-1})\right),\qquad
 U_e(s,t)=e^{i\eta(sB_e+tC_e)}U_e,
\end{equation}
and define $\Delta_U$ by the mixed derivative of this action.
There is no coupling prefactor in \eqref{nw-action}. At $U=1$,
$\Delta_1=d^Td$, where $d$ is the ordinary physical plaquette curl.

Let $U_e=e^{i\eta a_e}$, with
\begin{equation}\label{nw-sizes}
 \epsilon_0=\norm{a}_\infty,\qquad
 \epsilon_1=\norm{\nabla a}_\infty,\qquad
 F=\epsilon_1+\epsilon_0^2.
\end{equation}
The field can be complex; $\epsilon_0$ is bounded by a fixed small
constant. Comparing opposite plaquette sides by the integral formula
for matrix exponentials, and bounding the residual group commutator,
gives
\begin{equation}\label{nw-plaquette}
 \norm{P_p-I}+\norm{P_p^{-1}-I}\le C\eta^2F.
\end{equation}
Indeed opposite potentials differ by at most $\eta\epsilon_1$,
so their exponentials differ by $C\eta^2\epsilon_1$.
For the remaining two constant potentials the commutator costs
$C\eta^2\epsilon_0^2$. All path transports involved here have a
fixed number of links and norm at most $e^{C\eta\epsilon_0}$.

For the loop $e_1,\ldots,e_4$ around $p$, set
$V_i=U_{e_1}\cdots U_{e_{i-1}}$.
A negative traversal of the positive bond $f$ has exact chart
increment $B_{-f}=-\operatorname{Ad}_{U_f^{-1}}B_f$.
Put $X_i^B=\operatorname{Ad}_{V_i}B_{e_i}$ and
\begin{equation}\label{nw-curl}
 \mathcal C_U B=\eta^{-1}\sum_iX_i^B.
\end{equation}
The exact plaquette variation is
$P_p(s,t)=\prod_i e^{i\eta(sX_i^B+tX_i^C)}P_p$.
Every transport difference from identity is $O(\eta\epsilon_0)$,
and therefore
\begin{equation}\label{nw-curl-difference}
 \mathcal C_U=d+E_U,\qquad
 \norm{E_U}_{\infty\to\infty}
       +\norm{E_U^T}_{\infty\to\infty}\le C\epsilon_0.
\end{equation}
The transpose norm uses physical measure. Fixed plaquette incidence
also gives the same estimate in a weighted physical-cell Schur norm.

Writing $S_B=\sum_iX_i^B$, direct mixed differentiation gives
\begin{equation}\label{nw-mixed}
 \partial_s\partial_tP_p\big|_0=-\eta^2
 \left\{\frac{S_BS_C+S_CS_B}{2}
 +\frac12\sum_{i<j}
 ([X_i^B,X_j^C]+[X_i^C,X_j^B])\right\}P_p.
\end{equation}
For the oriented density $-\tau(P_p)$ the principal contribution is
$\eta^d\tau(((\mathcal C_UB)(\mathcal C_UC)
+(\mathcal C_UC)(\mathcal C_UB))P_p)/2$.
The remaining contribution is exactly
\begin{equation}\label{nw-commutator-remainder}
 \frac{\eta^d}{2\eta^2}\sum_{i<j}
 \tau\left(([X_i^B,X_j^C]+[X_i^C,X_j^B])(P_p-I)\right),
\end{equation}
because a commutator has trace zero. The reverse loop has curl
$-\operatorname{Ad}_{P_p^{-1}}\mathcal C_U$; cyclicity removes this
adjoint transport from its principal quadratic term. Thus
\begin{equation}\label{nw-hessian}
 \Delta_U=\mathcal C_U^TM_U\mathcal C_U+V_U,\qquad
 \tau(XM_{U,p}Y)=\frac14\tau((XY+YX)(P_p+P_p^{-1})),
\end{equation}
where $V_U$ is half \eqref{nw-commutator-remainder} plus half its
reverse-loop version. Projection onto the fixed traceless Lie space
is implicit in the definition of $M_U$. These coefficients are
bilinearly symmetric and holomorphic. Equation \eqref{nw-plaquette}
implies
\begin{equation}\label{nw-coefficients}
 \norm{M_U-I}\le C\eta^2F,
 \qquad \norm{V_U}_{\infty\to\infty}
       +\norm{V_U^T}_{\infty\to\infty}\le CF.
\end{equation}
Each curl has row and transpose-row norm $O(\eta^{-1})$.
Consequently the ordinary bounded local operator
$W_U=\Delta_U-\mathcal C_U^T\mathcal C_U$ has row and
transpose-row norm $CF$: the two inverse mesh factors are paid for
by $\eta^2$ in \eqref{nw-coefficients}. In particular
\begin{equation}\label{nw-small-form}
 \Delta_U-d^Td=d^TE_U+E_U^Td+E_U^TE_U+W_U
       =S_v^TV_{\rm kin}S_v,
 \quad S_v=(I,\nabla),
\end{equation}
with $\norm{V_{\rm kin}}_\kappa\le
C(\epsilon_0+\epsilon_1+\epsilon_0^2)$.
Here $\nabla$ contains all physical fine differences of bond
components; the curl is a fixed component combination of them.
The scalar gauge derivative is separately
\begin{equation}\label{nw-scalar-gradient}
 D_U\lambda(e)=\eta^{-1}
  (\operatorname{Ad}_{U_e}\lambda(e_+)-\lambda(e_-))
  =D\lambda+E_s\lambda,
 \qquad \norm{E_s}_\kappa\le C\epsilon_0.
\end{equation}
All coefficients in these two estimates have fixed fine range.

\subsection{The scalar compression and a form-resolvent lemma}
\label{nw-scalar}

For physical-cell block operators use the Banach algebra norm
\begin{equation}\label{nw-schur}
 \norm{T}_\kappa=\max\left\{
 \sup_x\sum_y e^{\kappa d(x,y)}\norm{T_{xy}},
 \sup_y\sum_x e^{\kappa d(x,y)}\norm{T_{xy}}\right\}.
\end{equation}
The norm is submultiplicative and invariant under transpose.
Fine-cell spaces have measure $\eta^d$, so it introduces no power
of the number of fine sites. The earlier flat inverse theorem gives
\begin{equation}\label{nw-flat-blocks}
 \norm{S_vG_0S_v^T}_\kappa\le K_v,
 \quad G_0=(d^Td+DR_0D^T+Q^Ta_{\rm pen}Q)^{-1},
\end{equation}
for $a_{\rm pen}$ in a fixed compact positive interval.
Both flat averages $Q=Q_M^{\rm st}$ and $Q_0$, and their transposes,
have uniform weighted norms by physical row/column contraction and
fixed final-scale range.

For clarity the additional scalar input is derived explicitly.
Let $G'_0=(\mathcal L+Q_0^Ta'Q_0)^{-1}$, $\mathcal L=D^TD$,
where $a'$ lies in a fixed compact positive interval, and let
$S_s=(I,D)$. In coarse Bloch aliases write $P_r$ for the fine
Laplacian multiplier, $a_r$ for the scalar average, and
\begin{align}
 A&=a_0a_0^-,&
 R_{s,k}&=\sum_{r\ne0}\frac{a_ra_r^-}{P_r^k},\nonumber\\
 D_{s,2}&=A+P_0^2R_{s,2},&
 N&=P_0+a'A+a'P_0R_{s,1}.                   \label{nw-scalar-denoms}
\end{align}
The scalar Woodbury formula has central entry
$(R_{s,1}+(a')^{-1})/[A+P_0(R_{s,1}+(a')^{-1})]$,
mixed entry $-a_0^-a_s/[P_s(A+P_0(R_{s,1}+(a')^{-1}))]$,
and a high diagonal $P_r^{-1}$ plus a rank-one correction.
These canceled expressions, high-alias gaps and alias Parseval
give $\norm{S_sG'_0S_s^T}_\kappa\le K_s$.
They also give the exact compression formula
\begin{equation}\label{nw-compression}
 C_0=Q_0(G'_0)^2Q_0^T=\frac{D_{s,2}}{N^2},
 \qquad C_0^{-1}=\frac{N^2}{D_{s,2}}.
\end{equation}
On an enlarged real Bloch cube, $A\ge c>0$ and the high remainders
are nonnegative and uniformly bounded. Thus $N,D_{s,2}$ are
bounded above and away from zero. Their uniformly bounded
holomorphic preliminary tube can be shrunk by a fixed amount to
retain these lower bounds. Periodic physical-cell reconstruction
and contour shifting prove
$\norm{C_0}_\kappa+\norm{C_0^{-1}}_\kappa\le K_c$.
The zero fiber gives $C_0=(a')^{-2}$; for $M=1$ it gives
$(P_0+a')^{-2}$. No inverse of the unpenalized zero mode is used.

We will repeatedly use the following elementary form identity.
If $G=\mathcal A^{-1}$, $K=SGS^T$ has weighted norm at most $K_*$,
and $I$ is a component of $S$, then for $K_*\norm{V}_\kappa<1$,
\begin{equation}\label{nw-form-resolvent}
 (\mathcal A+S^TVS)^{-1}
 =G-GS^TV(I+KV)^{-1}SG,
\end{equation}
and the augmented block has norm at most
$K_* /(1-K_*\norm{V}_\kappa)$ and differs from $K$ by at most
$K_*^2\norm{V}_\kappa/(1-K_*\norm{V}_\kappa)$.
The Neumann series in \eqref{nw-schur} proves these assertions and
holomorphy. Multiplication verifies the inverse at each fixed
spacing on the finite or infinite lattice. If the real baseline
$\mathcal A$ is positive definite, the same smallness gives
relative-form coercivity for real symmetric perturbations.

\subsection{Conditional closure of the projected gauge inverse}
\label{nw-conditional}

\begin{theorem}\label{nw-energy-theorem}
Suppose, in one common weighted algebra, that
\begin{gather}
 D_U=D+E_s,\quad Q_s(U)=Q_0+F_s,\quad Q_v(U)=Q+F_v,
 \quad \Delta_U=d^Td+S_v^TV_{\rm kin}S_v,\nonumber\\
 \norm{E_s}_\kappa+\norm{F_s}_\kappa+
 \norm{F_v}_\kappa+\norm{V_{\rm kin}}_\kappa\le\varepsilon.
                                                        \label{nw-hypotheses}
\end{gather}
Define
\begin{align}
 G'_U&=(D_U^TD_U+Q_s(U)^Ta'Q_s(U))^{-1},\nonumber\\
 C_U&=Q_s(U)(G'_U)^2Q_s(U)^T,\nonumber\\
 R_U&=I-G'_UQ_s(U)^TC_U^{-1}Q_s(U)G'_U,\nonumber\\
 \mathcal A_U&=\Delta_U+D_UR_UD_U^T
                     +Q_v(U)^Ta_{\rm pen}Q_v(U).       \label{nw-full-operator}
\end{align}
For a sufficiently small uniform $\varepsilon_*>0$, if
$\varepsilon\le\varepsilon_*$ all these inverses exist and
\begin{equation}\label{nw-energy-bound}
 \norm{S_sG'_US_s^T}_\kappa+
 \norm{C_U^{-1}}_\kappa+\norm{R_U}_\kappa+
 \norm{S_v\mathcal A_U^{-1}S_v^T}_\kappa\le C.
\end{equation}
Each displayed block differs from its flat counterpart by at
most $C\varepsilon$. Constants are uniform in refinement and
compatible periodic volume; the full-lattice statement also holds.
Holomorphic coefficient families give holomorphic inverse families.
On real data $R_U$ is the orthogonal projection off
$\operatorname{ran}(G'_UQ_s(U)^T)$.
\end{theorem}

\begin{proof}
The scalar precision difference is
\[
 D^TE_s+E_s^TD+E_s^TE_s+
 Q_0^Ta'F_s+F_s^Ta'Q_0+F_s^Ta'F_s.
\]
It equals $S_s^TV_sS_s$ with $\norm{V_s}_\kappa\le C_s\varepsilon$.
The form identity controls the scalar energy block and its
$O(\varepsilon)$ difference when $K_sC_s\varepsilon\le1/2$.
Product expansion gives $\norm{C_U-C_0}_\kappa\le C'_s\varepsilon$.
Equation \eqref{nw-compression} and
\[
 C_U^{-1}=(I+C_0^{-1}(C_U-C_0))^{-1}C_0^{-1}
\]
then control its inverse if $K_cC'_s\varepsilon\le1/2$.
The formula for $R_U$ gives its bound and its
$O(\varepsilon)$ difference. Its projection assertion is the Gram
projection formula; for complex coefficients it is a bilinear
holomorphic continuation rather than a positivity claim.

The gauge difference is exactly
\begin{equation}\label{nw-gauge-difference}
 D_UR_UD_U^T-DR_0D^T
 =D(R_U-R_0)D^T+E_sR_UD^T+DR_UE_s^T+E_sR_UE_s^T.
\end{equation}
The divergence $D^T$ on bonds is a fixed shifted component
combination of $\nabla$. Keeping these derivatives in $S_v$,
\eqref{nw-gauge-difference} is $S_v^TV_gS_v$ with norm
$C_g\varepsilon$. The average-penalty difference
$Q^Ta_{\rm pen}F_v+F_v^Ta_{\rm pen}Q+F_v^Ta_{\rm pen}F_v$
is a bounded $I,I$ coefficient of size $C_a\varepsilon$.
Thus $\mathcal A_U=\mathcal A_0+S_v^TV_{\rm tot}S_v$,
with $\norm{V_{\rm tot}}_\kappa\le C_v\varepsilon$.
A second application of \eqref{nw-form-resolvent}, with
$K_vC_v\varepsilon\le1/2$, proves the result.
\end{proof}

The theorem controls $G_U$, $\nabla G_U$, $G_U\nabla^T$ and
$\nabla G_U\nabla^T$, with exponential coarse-cell locality.
It does not propagate one-sided second derivatives by this
argument: differentiating $G_0S_v^T$ twice would require an
unavailable total third-order baseline block.

\subsection{Penalty independence and the homogeneous Wilson conclusion}
\label{nw-conclusion}

The scalar penalty is unnecessary for identifying the projection.
Fix $\mathcal L_U=D_U^TD_U$ and $Q_s$, and let
$G_b=(\mathcal L_U+Q_s^TbQ_s)^{-1}$,
$E_b=G_bQ_s^T$. Whenever the inverse and Gram inverse exist,
\begin{equation}\label{nw-range}
 \operatorname{ran}E_b
   =\{v:\mathcal L_Uv\in\operatorname{ran}Q_s^T\}.
\end{equation}
One inclusion follows by applying $\mathcal L_U$ to $G_bQ_s^T$.
For the other, $\mathcal L_Uv=Q_s^Tw$ implies
$v=G_bQ_s^T(w+bQ_sv)$.
The range is therefore independent of the symmetric coarse penalty
$b$. Its Gram inverse makes it nondegenerate, and
$E_b(E_b^TE_b)^{-1}E_b^T$ is the unique bilinear projection
onto this range along its orthogonal complement. Thus $R_U$ is
penalty independent. The identity holds directly for complex data
whenever the stated inverses exist, and on real data it is the
ordinary orthogonal projection identity.

We may consequently construct $R_U$ using $a'=1$.
It agrees with a differently penalized source formula wherever
that formula is defined, without estimating an unprinted scalar
penalty sequence. An estimate for that other scalar inverse
$G'_U$ itself still requires its own normalization bounds.

\begin{corollary}\label{nw-wilson-corollary}
Fix odd $L\ge3$, $M=L^j$, and the homogeneous full or periodic
lattice. Let $Q_v(U)$ be the deposited product of corrected
one-step derivatives at the ordinary reference trajectory, as
in \eqref{nb-actual-product}, and let $Q_s(U)$ be the deposited
scalar transported average \eqref{nb-scalar-average}.
Define $\Delta_U$ by \eqref{nw-action}, $D_U$ by
\eqref{nw-scalar-gradient}, and $R_U$ by
\eqref{nw-full-operator} with auxiliary scalar penalty one.
For $a_{\rm pen}$ in a fixed compact positive interval, there are
$\rho,\kappa,C>0$, independent of $j$ and compatible volume, such that
\begin{equation}\label{nw-final}
 \norm{a}_\infty+\norm{\nabla a}_\infty<\rho
 \quad\Longrightarrow\quad
 \norm{(I,\nabla)
 (\Delta_U+D_UR_UD_U^T+Q_v(U)^Ta_{\rm pen}Q_v(U))^{-1}
 (I,\nabla)^T}_\kappa\le C.
\end{equation}
The inverse is holomorphic on this complex global $C^1$ ball,
and its augmented block differs from the flat block by at most
$C(\norm{a}_\infty+\norm{\nabla a}_\infty)$.
\end{corollary}

\begin{proof}
Equations \eqref{nb-product-difference},
\eqref{nb-scalar-difference}, and \eqref{nb-cell-Schur}
give $\norm{Q_v(U)-Q}_\kappa+\norm{Q_s(U)-Q_0}_\kappa
\le C\epsilon_0$.
Their mechanism is the ordinary-reference estimate
$\norm{\log V_t}\le Ch_t\epsilon_0$, a physical one-step
coefficient difference $O(h_t\epsilon_0)$, and the geometric
sum $\sum_{t<j}h_t\le1/(L-1)$; no derivative of a different
nonlinear composition is substituted.
Equations \eqref{nw-small-form} and \eqref{nw-scalar-gradient}
supply the remaining hypotheses of Theorem~\ref{nw-energy-theorem}.
\end{proof}

The source convention check has a positive, concrete part.
Raw I lines 453--467 normalize $\operatorname{tr}1=1$, count
unoriented plaquettes once and define
$\mathsf A(U)=\sum_p(1-\operatorname{Re}\operatorname{tr}P_p)$;
lines 472--481 place the coupling separately. Thus this bare action
is exactly \eqref{nw-action} in dimension four.
Line 15645 uses the left chart $\mathsf A(e^{i\eta A}U_0)$,
and line 77577 gives the color pairing $\operatorname{tr}AB$.
Lines 13583--13593 define the full source operator
$\Delta+D_URD_U^T+Q^TaQ$, with $a=1$ at unit scale.
The remaining identification is to check that the reused
$\mathsf A$ in the Hessian passage denotes this same bare action
with the same physical summation convention. The separate source
claim about $\Delta-D_U^TD_U$ also requires matching its precise
one-form covariant curl. These are deposited-transcription
references, not newly recovered original journal proofs.

Corollary~\ref{nw-wilson-corollary} is a theorem for its explicitly
defined Wilson operator and deposited homogeneous averages.
It does not establish nonflat pointwise weak norms, nonflat current
domains, mixed retained histories, arbitrary nontrivial reference
connections, or an interacting continuum construction.
